% GENERATED FILE. Edit canonical lessons, then run npm run generate:books.
% canonical-source-hash: fnv1a64:54d6a177b89a62cc
% canonical-lessons: FR-C293-mystere, FR-C293-elucider, FR-C293-prouver, FR-C293-alibi, FR-C293-soupcon

\chapter{Un mystère}
\label{ch:fr-un-mystere}
\hlchaptermodality{293}

\begin{chapteropening}
\textbf{By the end of this chapter:} \emph{I can talk about solving a case: a mystery, clearing it up, proving something, an alibi and a suspicion.}

Complete the last lesson of chapter 293: I can talk about solving a case: a mystery, clearing it up, proving something, an alibi and a suspicion.
\end{chapteropening}

\section[le mystère]{le mystère — the mystery}

\label{lesson:FR-C293-mystere}



\begin{quote}
Before the new one: say the French for the clue, then the French for the print.
\end{quote}

\subsection*{You'll want to know: le mystère}
\textbf{le mystère} — \textquotedblleft{}the mystery\textquotedblright{}.

\begin{quote}
Personne ne sait où est le tableau: c'est un mystère.
\end{quote}

\begin{quote}
\emph{Nobody knows where the painting is: it is a mystery.}
\end{quote}

The adjective is \emph{mystérieux}.

\begin{cousinweb}
\emph{mystère} is a learned borrowing from Latin \emph{mystērium}, from Greek \emph{must\'{\={e}}rion} \textquotedblleft{}a secret rite\textquotedblright{} (Lewis \& Short).
\end{cousinweb}

\begin{grammarlens}[title={What you've built}]
One more word for telling, explaining and describing at B1.
\end{grammarlens}

\subsection*{Your turn}
\begin{itemize}
\raggedright
  \item \emph{Say it:} \emph{le mystère}
  \item \emph{Say it:} \emph{le mystère}, once more
  \item \emph{Say it:} a sentence of your own with \emph{le mystère}
  \item \emph{Recall:} say the French for the clue, then the French for the print, then say \emph{le mystère} again
\end{itemize}

\begin{tcolorbox}[breakable,colback=teal!4,colframe=teal!35!black,title={Before you move on}]
What does le mystère mean? (\textquotedblleft{}The mystery\textquotedblright{}.) Say it once more.
\end{tcolorbox}
\section[élucider]{élucider — to solve, to clear up}

\label{lesson:FR-C293-elucider}



\begin{quote}
Before the new one: say the French for the print, then the French for the mystery.
\end{quote}

\subsection*{You'll want to know: élucider}
\textbf{élucider} — \textquotedblleft{}to solve, to clear up\textquotedblright{}.

\begin{quote}
La police n'a jamais élucidé ce crime.
\end{quote}

\begin{quote}
\emph{The police never solved this crime.}
\end{quote}

Formal; in speech, \emph{résoudre}.

\begin{cousinweb}
\emph{élucider} is a learned borrowing from Late Latin \emph{ēlūcidāre} \textquotedblleft{}to make clear\textquotedblright{} (EtymDB), built on \emph{lūx} \textquotedblleft{}light\textquotedblright{}.
\end{cousinweb}

\begin{grammarlens}[title={What you've built}]
One more word for telling, explaining and describing at B1.
\end{grammarlens}

\subsection*{Your turn}
\begin{itemize}
\raggedright
  \item \emph{Say it:} \emph{élucider}
  \item \emph{Say it:} \emph{élucider}, once more
  \item \emph{Say it:} a sentence of your own with \emph{élucider}
  \item \emph{Recall:} say the French for the print, then the French for the mystery, then say \emph{élucider} again
\end{itemize}

\begin{tcolorbox}[breakable,colback=teal!4,colframe=teal!35!black,title={Before you move on}]
What does élucider mean? (\textquotedblleft{}To solve, to clear up\textquotedblright{}.) Say it once more.
\end{tcolorbox}
\section[prouver]{prouver — to prove}

\label{lesson:FR-C293-prouver}



\begin{quote}
Before the new one: say the French for the mystery, then the French for to solve.
\end{quote}

\subsection*{You'll want to know: prouver}
\textbf{prouver} — \textquotedblleft{}to prove\textquotedblright{}.

\begin{quote}
Il ne peut pas prouver qu'il était chez lui.
\end{quote}

\begin{quote}
\emph{He cannot prove that he was at home.}
\end{quote}

\emph{Prouver que} + indicative.

\begin{cousinweb}
\emph{prouver} is inherited from Latin \emph{probāre} \textquotedblleft{}to test, to prove, to approve\textquotedblright{} (EtymDB).
\end{cousinweb}

\begin{grammarlens}[title={What you've built}]
One more word for telling, explaining and describing at B1.
\end{grammarlens}

\subsection*{Your turn}
\begin{itemize}
\raggedright
  \item \emph{Say it:} \emph{prouver}
  \item \emph{Say it:} \emph{prouver}, once more
  \item \emph{Say it:} a sentence of your own with \emph{prouver}
  \item \emph{Recall:} say the French for the mystery, then the French for to solve, then say \emph{prouver} again
\end{itemize}

\begin{tcolorbox}[breakable,colback=teal!4,colframe=teal!35!black,title={Before you move on}]
What does prouver mean? (\textquotedblleft{}To prove\textquotedblright{}.) Say it once more.
\end{tcolorbox}
\section[l'alibi]{l'alibi — the alibi}

\label{lesson:FR-C293-alibi}



\begin{quote}
Before the new one: say the French for to solve, then the French for to prove.
\end{quote}

\subsection*{You'll want to know: l'alibi}
\textbf{l'alibi} — \textquotedblleft{}the alibi\textquotedblright{}.

\begin{quote}
Son alibi est solide: il était au cinéma.
\end{quote}

\begin{quote}
\emph{His alibi is solid: he was at the cinema.}
\end{quote}

Also an excuse: \emph{un bon alibi pour ne pas venir}.

\begin{cousinweb}
\emph{alibi} was borrowed from Latin \emph{alibī} \textquotedblleft{}elsewhere\textquotedblright{} (EtymDB), built on \emph{alius} \textquotedblleft{}other\textquotedblright{}.
\end{cousinweb}

\begin{grammarlens}[title={What you've built}]
One more word for telling, explaining and describing at B1.
\end{grammarlens}

\subsection*{Your turn}
\begin{itemize}
\raggedright
  \item \emph{Say it:} \emph{l'alibi}
  \item \emph{Say it:} \emph{l'alibi}, once more
  \item \emph{Say it:} a sentence of your own with \emph{l'alibi}
  \item \emph{Recall:} say the French for to solve, then the French for to prove, then say \emph{l'alibi} again
\end{itemize}

\begin{tcolorbox}[breakable,colback=teal!4,colframe=teal!35!black,title={Before you move on}]
What does l'alibi mean? (\textquotedblleft{}The alibi\textquotedblright{}.) Say it once more.
\end{tcolorbox}
\section[le soupçon]{le soupçon — the suspicion; a hint}

\label{lesson:FR-C293-soupcon}



\begin{quote}
Before the new one: say the French for to prove, then the French for the alibi.
\end{quote}

\subsection*{You'll want to know: le soupçon}
\textbf{le soupçon} — \textquotedblleft{}the suspicion; a hint\textquotedblright{}.

\begin{quote}
La police a des soupçons sur le voisin.
\end{quote}

\begin{quote}
\emph{The police have suspicions about the neighbour.}
\end{quote}

\emph{Un soupçon de sel} — a hint of salt.

\begin{cousinweb}
\emph{soupçon} is inherited from Late Latin \emph{suspectiō} \textquotedblleft{}suspicion\textquotedblright{} (EtymDB), from \emph{suspicere} \textquotedblleft{}to look up at, to suspect\textquotedblright{}.
\end{cousinweb}

\begin{grammarlens}[title={What you've built}]
One more word for telling, explaining and describing at B1.
\end{grammarlens}

\subsection*{Your turn}
\begin{itemize}
\raggedright
  \item \emph{Say it:} \emph{le soupçon}
  \item \emph{Say it:} \emph{le soupçon}, once more
  \item \emph{Say it:} a sentence of your own with \emph{le soupçon}
  \item \emph{Recall:} say the French for to prove, then the French for the alibi, then say \emph{le soupçon} again
\end{itemize}

\begin{tcolorbox}[breakable,colback=teal!4,colframe=teal!35!black,title={Before you move on}]
What does le soupçon mean? (\textquotedblleft{}The suspicion; a hint\textquotedblright{}.) Say it once more.
\end{tcolorbox}
